\documentclass{beamer}
\usepackage{xeCJK}
\usepackage{fontspec}        % XeLaTeX 必备
\setmonofont{WenQuanYi Micro Hei Mono}    % 等宽字体
\setCJKmainfont{WenQuanYi Micro Hei}
\usepackage{xcolor}

\usetheme{Madrid}

\title{Python sys.monitoring 简介}
\subtitle{面向 Tracing / Debugging / Profiling 的新一代运行时观测接口}
\author{}
\date{}

\begin{document}

%------------------------------------------------
\begin{frame}
  \titlepage
\end{frame}

%------------------------------------------------
\begin{frame}{背景：为什么需要 sys.monitoring？}
\begin{itemize}
  \item sys.settrace 存在的问题：
    \begin{itemize}
      \item 事件粒度粗（行级）
      \item 全局单一回调，难以组合多个工具
      \item 性能开销大
    \end{itemize}
  \item Python 3.12 引入 sys.monitoring
  \item 目标：
    \begin{itemize}
      \item 更细粒度事件
      \item 多工具并存（debugger / profiler / tracer）
      \item 更低运行时开销
    \end{itemize}
\end{itemize}
\end{frame}

%------------------------------------------------
\begin{frame}{sys.monitoring 核心机制}
\begin{itemize}
  \item 事件驱动模型（Event-based）
  \item 为不同工具分配 \textbf{tool\_id}
  \item 每个工具订阅自己关心的事件
\end{itemize}

\textbf{典型事件类型：}
\begin{itemize}
  \item PY\_START / PY\_RETURN
  \item LINE
  \item CALL / C\_CALL
  \item INSTRUCTION
\end{itemize}
\end{frame}

%------------------------------------------------
\begin{frame}[fragile]{Tracer 示例：函数调用追踪}
\begin{verbatim}
import sys
from sys import monitoring

TOOL_ID = 1

def on_call(code, offset):
    print(f"CALL {code.co_name}")

monitoring.use_tool_id(TOOL_ID, "simple-tracer")
monitoring.register_callback(TOOL_ID,
    monitoring.events.PY_CALL, on_call)
monitoring.set_events(TOOL_ID, monitoring.events.PY_CALL)

def foo():
    return 42
foo()
\end{verbatim}
\end{frame}

%------------------------------------------------
\begin{frame}[fragile]{Debugger / Profiler 场景}
\textbf{Debugger：}
\begin{itemize}
  \item 监听 LINE / INSTRUCTION 事件
  \item 支持断点、单步执行
\end{itemize}

\textbf{Profiler：}
\begin{itemize}
  \item 监听 CALL / RETURN
  \item 构建调用栈、统计时间
\end{itemize}

\textbf{Tracer 示例代码：}
\begin{verbatim}
# 调试器和性能分析器可以同时启用
# 而不会覆盖彼此的回调逻辑
\end{verbatim}
\end{frame}

%------------------------------------------------
\begin{frame}[fragile]{sys.monitoring vs sys.settrace}
\begin{tabular}{l|c|c}
维度 & sys.settrace & sys.monitoring \\
\hline
事件粒度 & 行级 & 行 / 指令 / 调用 \\
多工具支持 & 否 & 是（tool\_id） \\
性能 & 较高开销 & 更低 \\
设计年代 & Python 早期 & Python 3.12+ \\
适用场景 & 简单调试 & 可观测性 / 运行时分析 \\
\end{tabular}

\begin{verbatim}
# sys.monitoring 支持更高性能的 tracing
# 可同时用于调试和性能分析
\end{verbatim}
\end{frame}

%------------------------------------------------
\begin{frame}[fragile]{总结}
\begin{itemize}
  \item sys.monitoring 是 Python 运行时可观测性的基础设施
  \item 更适合：
    \begin{itemize}
      \item 高性能 tracing
      \item 组合式调试与分析工具
      \item 系统级行为建模
    \end{itemize}
  \item 可作为：
    \begin{itemize}
      \item 动态分析工具
      \item 运行时剖析框架
      \item 可观测性驱动优化的底层支撑
    \end{itemize}
\end{itemize}

\begin{verbatim}
# 使用 fragile + verbatim, Python f-string 安全显示
\end{verbatim}
\end{frame}

\end{document}
